\section{Exact sourced networks and convergent rational integration}
\label{networks:section}

We work in the purely holomorphic case with two-dimensional source
spinors.  The graph, the order and orientation of its edges, and every
original source remain fixed throughout a reduction.  This section proves
an exact evaluation by a convergent rational series at every finite current
order.  A finite expression in a prescribed set of periods is a separate
property, considered in Theorem~\ref{periods:finite-span}.

\subsection{Original currents, projective measure, and signs}
\label{networks:conventions}

Let $\Gamma$ be a connected ordered graph with $N$ edges, loop number
$n=N-|V|+1\ge1$, and $N=2n+1$.  Thus $|V|=n+2$.  Its source spinors obey
$\sum_{v\in V}\lambda_{v,+\dot a}=0$.  The two dotted indices are
oriented by $\epsilon^{\dot+\dot-}=1$, and we write
\begin{equation}
 \langle\lambda_u,\lambda_v\rangle
 =\epsilon^{\dot a\dot b}\lambda_{u,+\dot a}\lambda_{v,+\dot b},
 \qquad
 d^2y_e=\frac12\epsilon^{\dot a\dot b}
 dy_{e,+\dot a}\wedge dy_{e,+\dot b}.
 \label{networks:spinor-volume}
\end{equation}
The brackets contract only the displayed indices.  The indices of an
inserted current word remain free.

Choose one ground vertex and delete its column from the incidence matrix.
For an edge $e$ directed from $e(0)$ to $e(1)$, its row is
$B_{ev}=\delta_{v,e(0)}-\delta_{v,e(1)}$.  Positive Schwinger parameters
$t_e$ are resistances.  Kirchhoff's equations and their unique grounded
solution are
\begin{equation}
 B^{\mathsf T}y_{+\dot a}=\lambda_{+\dot a},\qquad
 t\,y_{+\dot a}=B\varphi_{+\dot a},\qquad
 \Lambda=B^{\mathsf T}t^{-1}B,\qquad
 \varphi_{+\dot a}=\Lambda^{-1}\lambda_{+\dot a}.
 \label{networks:kirchhoff}
\end{equation}
Here $t$ is the diagonal resistance matrix.  Its positive quadratic form
makes $\Lambda$ invertible.

For a spanning tree $T$, order its complementary edges as
$c_1<\cdots<c_n$.  The fundamental cycle matrix $C_T$ has entry $+1$ on
its own complementary edge and the signed return path on the tree.
Define $M=C_TtC_T^{\mathsf T}$ and $U_\Gamma=\det M$.  Cauchy--Binet,
together with the unimodularity of the cycle minors, gives
\begin{equation}
 U_\Gamma(t)=\sum_{T'\text{ spanning tree}}
                  \prod_{e\notin T'}t_e.
 \label{networks:kirchhoff-polynomial}
\end{equation}
Let $y^T$ be the unique source flow supported on $T$.  Every original edge
current has the exact affine expression
\begin{equation}
 y_{e,+\dot a}=y^T_{e,+\dot a}
       +\sum_{r=1}^n(C_T)_{re}y_{c_r,+\dot a}.
 \label{networks:affine-original-currents}
\end{equation}
This identity is used before integrating a current polynomial; in
particular it retains constant source terms in words of mixed degree.

The cycle forms are
$\psi_{r,+\dot a}=\sum_e(C_T)_{re}y_{e,+\dot a}\,dt_e$.
The kernel convention used here is
\begin{equation}
 \Omega_\Gamma=U_\Gamma^{-2}
       \bigwedge_{\dot a=\dot+,\dot-}\ \bigwedge_{r=1}^n\psi_{r,+\dot a}
       =D_\Gamma(t)\,\iota_{\sum_e t_e\partial_{t_e}}
                          (dt_0\wedge\cdots\wedge dt_{N-1}).
 \label{networks:area-definition}
\end{equation}
Differentiating $C_Tty=0$ at fixed original sources and using
\eqref{networks:affine-original-currents} gives
$\psi_{r,+\dot a}=-\sum_sM_{rs}dy_{c_s,+\dot a}$.
The two determinants cancel $U_\Gamma^{-2}$; interleaving their cycle rows
gives the exact current-area identity
\begin{equation}
 \boxed{\Omega_\Gamma=(-1)^{\binom n2}
              \bigwedge_{r=1}^n d^2y_{c_r}.}
 \label{networks:current-area}
\end{equation}
A unimodular change of cycle basis has squared determinant one, so this
formula does not depend on the spanning tree.  The positive projective
orientation is the radial contraction in
\eqref{networks:area-definition}.  For any gauge edge $g$, the equivalent
cofactor formula is
\begin{equation}
 D_\Gamma(t)=\frac{(-1)^{g+\binom n2}}{t_g}
 \det\left(\frac{\partial y_{c_r,+\dot a}}{\partial t_f}\right)
 _{(r,\dot a),\ f\ne g}.
 \label{networks:density-cofactor}
\end{equation}
Columns retain the original edge order.  Reversing an edge leaves the
bare two-dimensional kernel unchanged and negates that edge's inserted
current.  Permuting the edges multiplies the kernel by the parity of the
permutation.  These signs are never absorbed into a new edge convention.
For comparison with the current-area convention of~\cite{li2026},
$\mathcal I_\Gamma=(-1)^{n(n+1)/2}K_\Gamma$.

The complete current moment with free indices is
\begin{equation}
 K_\Gamma[e_1,\ldots,e_m]_{+\dot a_1\cdots+\dot a_m}
 =\int_{\mathbb{RP}^{N-1}_{>0}}
       \prod_{j=1}^m y_{e_j,+\dot a_j}\,\Omega_\Gamma.
 \label{networks:moment}
\end{equation}
It is the coefficient obtained by applying the corresponding $m$ operators
$-\partial/\partial z_{e_j}^{+\dot a_j}$ to
$\int e^{-\sum_e y_{e,+\dot a}z_e^{+\dot a}}\Omega_\Gamma$ at $z=0$.
All statements below also apply to finite linear combinations of these
words, including different degrees.

\subsection{Source-preserving elimination on every product forest}
\label{networks:elimination}

A marked product forest has ordered distinguished vertices $A,B$, the
marked edge $AB$, and $n$ interaction vertices.  Removing $A,B$ leaves
trees.  A component with $k$ interaction vertices has $k-1$ internal edges
and $k+1$ spokes, and meets both marked vertices.  These are the graph
conditions arising in the product-forest classification of~\cite{li2026}.
The following network argument permits arbitrary balanced sources at all
vertices; it does not impose a cubic source-support restriction.

\begin{proposition}[Terminating sourced elimination]
\label{networks:kron-theorem}
Choose a root in each interaction tree and remove descendants before
their parents.  Every removed vertex has at most three retained
neighbors.  Exact Schur complementation terminates at the marked resistor,
preserves all original source coefficients, and reconstructs every
original current by reverse substitution.
\end{proposition}
\begin{proof}
After the descendants of $v$ are removed, its only possible retained
neighbors are its parent and $A,B$.  Eliminating a vertex adds edges only
between its retained neighbors.  Such an operation cannot connect two
distinct unprocessed interaction branches.  This proves the assertion
inductively, including a root of arbitrarily high original valence.

Write $g_{vu}=1/t_{vu}$ for the present conductance, allowing zero for an
absent edge, and $g_v=\sum_u g_{vu}$.  The equation at $v$ gives
\begin{equation}
 \varphi_{v,+\dot a}
 =\frac{\lambda_{v,+\dot a}}{g_v}
       +\sum_u\frac{g_{vu}}{g_v}\varphi_{u,+\dot a}.
 \label{networks:reverse-potential}
\end{equation}
Substitution into the retained equations yields
\begin{equation}
 \boxed{g^{\rm red}_{uw}=g_{uw}+\frac{g_{vu}g_{vw}}{g_v},\qquad
 \lambda^{\rm red}_{u,+\dot a}
       =\lambda_{u,+\dot a}+\frac{g_{vu}}{g_v}\lambda_{v,+\dot a}.}
 \label{networks:sourced-schur}
\end{equation}
All pivots are positive, and the transfer coefficients sum to one.
Consequently the complete source conservation equation survives each
step.  At this stage $\lambda_v$ means the source accumulated from every
previously eliminated vertex, with its full resistance dependence.
At the end, ground $B$ and use
$\varphi_A=\lambda_A^{\rm final}/g_{AB}^{\rm final}$.
Applying \eqref{networks:reverse-potential} in reverse order recovers all
potentials and hence $y_{e,+\dot a}=(\varphi_{e(0),+\dot a}
-\varphi_{e(1),+\dot a})/t_e$ on every original edge.
\end{proof}

For example, the complete affine current map of one eliminated star can
be displayed without solving the rest of the network.  Orient its edges
outward from $v$, and orient a generated triangle edge from $u$ to $w$.
Let $y^{(v)}_{uw}$ denote the current in just the generated parallel branch.
Then
\begin{align}
 y^{(v)}_{uw,+\dot a}
 &=\frac{g_{vu}g_{vw}/g_v}{g^{\rm red}_{uw}}y^{\rm red}_{uw,+\dot a},
 &y^{\rm old}_{uw,+\dot a}
 &=\frac{g_{uw}}{g^{\rm red}_{uw}}y^{\rm red}_{uw,+\dot a},\\
 y_{vu,+\dot a}
 &=\frac{g_{vu}}{g_v}\lambda_{v,+\dot a}
       -\sum_{w\ne u}y^{(v)}_{uw,+\dot a}.
 \label{networks:full-star-current}
\end{align}
Antisymmetry gives $y^{(v)}_{wu}=-y^{(v)}_{uw}$, so summing the last line
recovers $\sum_u y_{vu}=\lambda_v$.  Every old parallel branch and every
source offset is separately retained.

The determinant part of the Schur complement also gives an unexpanded
Kirchhoff polynomial:
\begin{equation}
 U_\Gamma(t)
 =\left(\prod_{e\in E}t_e\right)g_{AB}^{\rm final}
       \prod_{v\text{ eliminated}}g_v.
 \label{networks:pivot-U}
\end{equation}
Indeed $U_\Gamma=(\prod_e t_e)\det\Lambda$, and each elimination removes
one positive pivot.  The final grounded determinant is
$g_{AB}^{\rm final}$.  There are linearly many local steps and at most
quadratically many arithmetic operations on complete source-coefficient
rows.  This is not a bound on rational bit sizes.

Intermediate networks need not satisfy $N=2n+1$.  Formula
\eqref{networks:sourced-schur} remains valid on them, while the area form
is always the original form \eqref{networks:current-area}.  In particular,
one cannot substitute a lower-loop tight-graph kernel for a nontight
intermediate network.

\subsection{All differentials of transferred sources}
\label{networks:differentials}

The transferred source depends on the resistances.  Differentiating
\eqref{networks:sourced-schur} includes the term
$\lambda_{v,+\dot a}\,d(g_{vu}/g_v)$ as well as the differentials of
sources accumulated earlier.  Freezing a transferred source changes the
area form.  A convenient exact evaluation differentiates the reconstructed
original network instead.  At fixed original sources,
\begin{equation}
 \frac{\partial\varphi_{+\dot a}}{\partial t_f}
 =\Lambda^{-1}B_f^{\mathsf T}\frac{y_{f,+\dot a}}{t_f},\qquad
 \boxed{
 \frac{\partial y_{e,+\dot a}}{\partial t_f}
 =\left(\frac{B_e\Lambda^{-1}B_f^{\mathsf T}}{t_e}
                  -\delta_{ef}\right)\frac{y_{f,+\dot a}}{t_f}.}
 \label{networks:full-derivative}
\end{equation}
This follows from
$\partial_{t_f}\Lambda=-B_f^{\mathsf T}B_f/t_f^2$ and
\eqref{networks:kirchhoff}.  Equations
\eqref{networks:full-derivative} and \eqref{networks:density-cofactor}
therefore include every source-transfer differential without separately
expanding them.

An important local instance is a sourced degree-two vertex $v$ between
$u,w$, with a direct edge $uw$.  Orient the three original currents as
$u\to v,v\to w,u\to w$.  If the retained resistance and current are
$T$ and $J$, positive coordinates $0<r,z<1$ give
\begin{align}
 t_{uv}=rT/z,\quad t_{vw}=(1-r)T/z,\quad t_{uw}=T/(1-z),\\
 \lambda_u^{\rm red}=\lambda_u+(1-r)\lambda_v,\qquad
 \lambda_w^{\rm red}=\lambda_w+r\lambda_v,\\
 y_{uv,+\dot a}=zJ_{+\dot a}-(1-r)\lambda_{v,+\dot a},\\
 y_{vw,+\dot a}=zJ_{+\dot a}+r\lambda_{v,+\dot a},\\
 y_{uw,+\dot a}=(1-z)J_{+\dot a}.
 \label{networks:ear-map}
\end{align}
Choose a reduced spanning tree containing $uw$.  Modulo its other chord
differentials, source conservation gives $dJ=-\lambda_v\,dr$.
The new chord $(1-z)J$ has area
$(1-z)\langle\lambda_v,J\rangle\,dr\wedge dz$.
The local resistance Jacobian in order $(uv,vw,uw)$ is
$-T^2/[z^3(1-z)^2]$.  Combining its sign with the ratio of the signs in
\eqref{networks:current-area} gives the complete density factor
\begin{equation}
 (-1)^n\operatorname{sgn}(\text{original edge permutation})
 (1-z)\langle\lambda_v,J\rangle\,dr\,dz.
 \label{networks:ear-density}
\end{equation}
The reduced kernel is evaluated at the full sources in
\eqref{networks:ear-map}; its retained coordinates are differentiated at
fixed $r$.  The displayed local factor already accounts for the varying
source in the $r$ direction.  A current word is expanded using all three
affine maps.  Each splitting monomial integrates exactly as
\begin{equation}
 \int_0^1 z^a(1-z)^{b+1}dz=\frac{a!(b+1)!}{(a+b+2)!}.
 \label{networks:ear-beta}
\end{equation}
The remaining source transfer is polynomial at any finite current order
and is integrated by the ordinary interval moments.  This local formula
requires the removed source to be fixed during that step; the general
rule \eqref{networks:full-derivative} applies without that restriction.

\subsection{Series chains and complete scalar reduction}
\label{networks:chains}

Consider a maximal path from a retained vertex $p$ to $q$.  Orient it
consistently and let $s_j=\pm1$ compare that direction with the original
direction of its $j$th edge.  Write its resistance as $t_j=Rx_j$,
$\sum_{j=1}^Lx_j=1$, and denote its internal vertices by
$u_1,\ldots,u_{L-1}$.  Define the complete cumulative sources and their
weighted average by
\begin{equation}
 Q_{j,+\dot a}=\sum_{k<j}\lambda_{u_k,+\dot a},\qquad
 \widehat\lambda_{+\dot a}
   =\sum_{k=1}^{L-1}\left(\sum_{j>k}x_j\right)\lambda_{u_k,+\dot a}.
 \label{networks:chain-sources}
\end{equation}
If $Y$ is the current of the retained path, every original current and
both retained source rows are
\begin{align}
 y_{j,+\dot a}&=s_j\bigl(Y_{+\dot a}-\widehat\lambda_{+\dot a}
                                  +Q_{j,+\dot a}\bigr),\\
 \lambda_{p,+\dot a}^{\rm red}&=\lambda_{p,+\dot a}
                                      +\widehat\lambda_{+\dot a},\\
 \lambda_{q,+\dot a}^{\rm red}&=\lambda_{q,+\dot a}
              +\sum_k\lambda_{u_k,+\dot a}-\widehat\lambda_{+\dot a}.
 \label{networks:chain-map}
\end{align}
Conservation along the path gives successive differences $\lambda_{u_k}$,
and its voltage drop is $R Y=\sum_jt_js_jy_j$.  These two equations prove
\eqref{networks:chain-map}.  Thus neither endpoint source nor the offsets
of any removed edge are omitted.

The cycle columns on a consistently oriented series path coincide.
Consequently, for positive Lebesgue measures,
\begin{equation}
 U_\Gamma(t)=U_{\rm core}(R),\qquad
 \prod_e dt_e=\prod_{\text{paths}}R^{L-1}\,dR\,d^{L-1}x.
 \label{networks:chain-measure}
\end{equation}
For a complete numerator monomial with nonnegative exponents $k_e$, the
exact push-forward on one path is
\begin{equation}
 \boxed{\alpha_{\rm path}=\sum_{e\in\rm path}(k_e+1),\qquad
 \int_{\Delta^{L-1}}\prod_{e\in\rm path}x_e^{k_e}d^{L-1}x
       =\frac{\prod_{e\in\rm path}k_e!}{(\alpha_{\rm path}-1)!}.}
 \label{networks:chain-factor}
\end{equation}
This changes neither the loop number nor the denominator exponent.

For completeness, the spanning-tree current formula is
\begin{equation}
 U_\Gamma y_{e,+\dot a}
 =\sum_{T'}\left(\prod_{f\notin T'}t_f\right)y^{T'}_{e,+\dot a}.
 \label{networks:tree-current}
\end{equation}
It follows from Cramer's rule for \eqref{networks:kirchhoff} and the same
tree-minor expansion as \eqref{networks:kirchhoff-polynomial}.  Thus the
current is a convex weighted average of finitely many source flows and
is bounded on positive times.

The bare Schwinger numerator $P=U_\Gamma^{n+2}D_\Gamma$ is polynomial
of time degree $n^2-1$.  Start with the $\psi$-row representation
\eqref{networks:area-definition}, before passing to the current Jacobian.
Clearing one factor $U_\Gamma$ from each of its $2n$ rows gives denominator
$t_gU_\Gamma^{2n+2}$ and polynomial entries
$(C_T)_{re}(U_\Gamma y_{e,+\dot a})$, with column $g$ omitted.
At a simple zero of any irreducible factor
of the multiaffine $U_\Gamma$, the cycle matrix $M$ has corank one.
Its adjugate makes the two source blocks proportional there; subtracting
one block from the other supplies at least $n$ factors of that
irreducible polynomial.  Changing the omitted edge in
\eqref{networks:density-cofactor} removes the apparent factor $1/t_g$.
Since a multiaffine polynomial has no repeated nonconstant factor, the
division is exact and the $n$ cancelled factors leave denominator
$U_\Gamma^{n+2}$.  Homogeneity then gives the asserted degree.
A word of degree $m$ therefore gives the complete scalar reduction
\begin{equation}
 \frac{P(t)\prod_{j=1}^m(U_\Gamma y_{e_j,+\dot a_j})}
      {U_\Gamma^{q}},\qquad
 q=n+2+m,\qquad
 \deg_t\bigl(P\prod_j U_\Gamma y_{e_j}\bigr)=n^2-1+nm.
 \label{networks:scalar-numerator}
\end{equation}
Its source degree is $2n+m$, with all free indices retained.  Apply
\eqref{networks:chain-factor} to every collected monomial, including its
sign and full source coefficient.  The strict convergence criterion of
Proposition~\ref{periods:integer-convergence} applies to these monomials.
The compact construction below independently proves absolute convergence
of the complete moments.

\subsection{A compact cube atlas with a positive denominator}
\label{networks:cubes}

If a bridge is present, its resistance changes no current; together with
common scaling this makes the current Jacobian deficient and all moments
\eqref{networks:moment} vanish.  Assume henceforth that the graph is
bridgeless.  Let $\tau_\Gamma$ be its number of spanning trees.
For each spanning tree $T$, choose an ordering $c_1,\ldots,c_n$ of its
complementary edges.  The ordering here is by increasing resistance,
not necessarily the original edge order.  Put
\begin{equation}
 t_{c_n}=1,\qquad
 t_{c_j}=\prod_{k=j}^{n-1}r_k\quad(j<n),\qquad
 t_f=x_f t_{c_{j(f)}}\quad(f\in T),
 \label{networks:cube-map}
\end{equation}
where $j(f)$ is the first complementary edge whose fundamental cycle
contains $f$.  All $x_f,r_k$ range independently over $(0,1)$.
There are $|T|+n-1=N-1=2n$ coordinates.

\begin{proposition}[Complete compact atlas]
\label{networks:atlas}
The $\tau_\Gamma n!$ maps \eqref{networks:cube-map} cover the positive
projective resistance space exactly once away from equality walls.
After division by the selected complementary monomial, the Kirchhoff
polynomial becomes a polynomial $\widetilde U$ with nonnegative integer
coefficients and constant term one.  On every closed cube,
\begin{equation}
              1\le\widetilde U\le\tau_\Gamma.
 \label{networks:cube-bounds}
\end{equation}
Every original current extends as a rational function without poles to
the closed cube.
\end{proposition}
\begin{proof}
The minimum-spanning-tree cycle criterion is
$t_f\le t_c$ whenever $f$ belongs to the tree path of the chord $c$.
With the complementary order fixed, these are exactly the inequalities
in \eqref{networks:cube-map}.  The minimum tree and complementary order
are unique away from ties, proving the covering assertion.

For a prefix $c_1,\ldots,c_k$, include those tree edges whose upper
resistance in \eqref{networks:cube-map} belongs to the prefix.  The
resulting subgraph has cycle rank $k$: it contains the fundamental path
of each prefix chord, while its tree edges are acyclic.  Every other
spanning-tree complement contains at least $k$ edges in this subgraph.
These inequalities say precisely that each exponent of a successive
ratio $r_k$ in
$\prod_{e\notin T'}t_e/\prod_{e\notin T}t_e$ is nonnegative.
Tree-fraction exponents are already zero or one.  The chosen complement
gives the unique constant monomial.  At the all-one corner each tree
monomial is one, proving \eqref{networks:cube-bounds}.
Divide \eqref{networks:tree-current} by the same monomial.  Every numerator
is now polynomial and the common denominator is $\widetilde U\ge1$.
\end{proof}

Order the cube coordinates by the original tree-edge order, followed by
$r_1,\ldots,r_{n-1}$.  If $t_e=\prod_jx_j^{a_{ej}}$ in these coordinates,
the projective orientation sign is
\begin{equation}
 \chi_{T,c}=\det\bigl[\mathbf1\mid(a_{ej})\bigr]\in\{1,-1\},\qquad
 \det[t,\partial_1t,\ldots,\partial_{2n}t]
       =\chi_{T,c}\frac{\prod_e t_e}{\prod_jx_j}.
 \label{networks:cube-sign}
\end{equation}
The determinant has absolute value one because the inverse coordinates
are the successive ratios and tree fractions in
\eqref{networks:cube-map}.  Permuting whole two-row chord blocks has even
parity.  Hence the complete signed cube integrand is
\begin{equation}
 \chi_{T,c}(-1)^{\binom n2}
 \det\left(\frac{\partial y_{c_r,+\dot a}}{\partial x_j}\right)
 \prod_{\ell=1}^m y_{e_\ell,+\dot a_\ell}.
 \label{networks:cube-integrand}
\end{equation}
It is continuous on the closed cube.  The finite atlas proves absolute
integrability of every finite original-current moment, including arbitrary
balanced rational source assignments.

\subsection{Polynomial coefficients and an explicit rational tail}
\label{networks:series}

On one cube write the full current numerators as
$Y_{e,+\dot a}=\widetilde U y_{e,+\dot a}$.
Arrange all chord spinors in their complete cycle order.  The augmented
polynomial determinant satisfies
\begin{equation}
 \boxed{
 \det(\partial_j y_{c_r,+\dot a})
 =\frac{\det\begin{pmatrix}
 \widetilde U&\partial_1\widetilde U&\cdots&\partial_{2n}\widetilde U\\
 Y&\partial_1Y&\cdots&\partial_{2n}Y
 \end{pmatrix}}{\widetilde U^{2n+1}}.}
 \label{networks:augmented}
\end{equation}
Subtract $y_{c_r,+\dot a}$ times the first row from its corresponding
row.  The first column then has only its top entry, and the other entries
are $\widetilde U\partial_jy_{c_r,+\dot a}$, proving the identity.
All derivatives hold the original sources fixed.

A degree-$m$ current word initially gives denominator exponent
$2n+1+m$.  Exact polynomial division may remove factors of
$\widetilde U$; each zero remainder must be checked.  Group terms of mixed
degree without discarding any term, and combine complete numerators with
identical denominators and powers.  Label the retained components by $c$,
with full numerator $P_c$ and positive denominator exponent $\nu_c$.
Components with exponent zero are integrated as polynomials separately.

\begin{theorem}[Complete rational series with a certified finite tail]
\label{networks:rational-series}
Every finite current moment on a product forest has a finite specification
as a convergent series of rational coefficients at rational sources.
Each coefficient is computed by finite polynomial multiplication and
monomial integration.  For every positive rational error tolerance a
finite sufficient order is determined by an explicit rational bound.
The assertion holds coefficientwise for the complete source polynomial
and every free spinor index.
\end{theorem}
\begin{proof}
Set the expansion center to $(\tau_\Gamma+1)/2$.  By
\eqref{networks:cube-bounds}, the absolute expansion ratio is bounded by
$(\tau_\Gamma-1)/(\tau_\Gamma+1)<1$.  The exact binomial expansion is
\begin{equation}
 \frac{P_c}{\widetilde U_c^{\nu_c}}
 =\sum_{k=0}^{\infty}
 \frac{\binom{\nu_c+k-1}{k}}
      {((\tau_\Gamma+1)/2)^{\nu_c+k}}
 P_c\bigl((\tau_\Gamma+1)/2-\widetilde U_c\bigr)^k.
 \label{networks:binomial-series}
\end{equation}
It converges absolutely and uniformly on the closed cube.  If the
complete polynomial in its $k$th summand has coefficients $p_{c,k,\beta}$,
its integral is the finite rational number
\begin{equation}
 \sum_{\beta}p_{c,k,\beta}
           \prod_{j=1}^{2n}\frac1{\beta_j+1},\qquad
 \int_{[0,1]^{2n}}\prod_jx_j^{\beta_j}d^{2n}x
                   =\prod_j\frac1{\beta_j+1}.
 \label{networks:cube-moments}
\end{equation}
All charts and all polynomial terms contribute to each complete level.

Write $P_c=\sum_\beta p_{c,\beta}x^\beta$ and define its rational absolute
moment by
$A_c=\sum_\beta|p_{c,\beta}|\prod_j(\beta_j+1)^{-1}$.
After order $K$, an absolute bound for the entire remainder is
\begin{equation}
 \boxed{
 B_K=\left(\frac{\tau_\Gamma-1}{\tau_\Gamma+1}\right)^{K+1}
 \sum_c A_c\sum_{j=0}^{\nu_c-1}
       \binom{K+j}{j}\left(\frac2{\tau_\Gamma+1}\right)^j.}
 \label{networks:tail}
\end{equation}
For a positive integer $\nu$ and $0\le r<1$, the identity behind this
formula is
\begin{equation}
 (1-r)^\nu\sum_{k=K+1}^{\infty}\binom{\nu+k-1}{k}r^k
 =r^{K+1}\sum_{j=0}^{\nu-1}\binom{K+j}{j}(1-r)^j.
 \label{networks:finite-tail-identity}
\end{equation}
For $\nu=1$ this is the geometric tail.  Pascal's identity proves the
general case by induction; equivalently, both sides have identical
successive differences in $K$ and tend to zero.  Apply it with
$1-r=2/(\tau_\Gamma+1)$ and use
$\int|P_c|\le A_c$.
The bound is decreasing and tends to zero, being a geometric factor
times a polynomial of finite degree in $K$.  If a width $\eta>0$ is
requested, the first $K$ with $2B_K\le\eta$ is a sufficient order before
any of the corresponding series coefficients are computed.  The exact
partial sum plus or minus $B_K$ is a certified rational interval.
\end{proof}

At word degree $m$, the numerator has source degree $2n+m$.  Hence there
are finitely many source coefficients, each covered by the theorem.
For a fixed rational source assignment the same construction retains the
entire free-index tensor.  A finite specification consists of the original
graph and sources, the complete component polynomials, their powers, and
\eqref{networks:binomial-series}--\eqref{networks:tail}; it contains no
unevaluated resistance integral.  A finite partial sum alone is not that
specification.  Storage or order limits may halt a computation, but must
not turn a partial result into an exact answer.

\subsection{Two complete analytic controls}
\label{networks:controls}

\paragraph{The triangle, including its full rational series.}
Take ordered edges $(01,02,12)$ and arbitrary balanced sources
$\lambda_2=-\lambda_0-\lambda_1$.  In the simplex gauge
$t_{01}+t_{02}+t_{12}=1$, put $x=t_{02}$ and $y=t_{12}$.  Direct solution
of \eqref{networks:kirchhoff} gives
\begin{align}
 y_{01,+\dot a}&=x\lambda_{0,+\dot a}-y\lambda_{1,+\dot a},\\
 y_{02,+\dot a}&=(1-x)\lambda_{0,+\dot a}+y\lambda_{1,+\dot a},\\
 y_{12,+\dot a}&=x\lambda_{0,+\dot a}+(1-y)\lambda_{1,+\dot a},\\
 \Omega_\Gamma&=-\langle\lambda_0,\lambda_1\rangle\,dx\,dy,
 \qquad x,y>0,\quad x+y<1.
 \label{networks:triangle-chart}
\end{align}
The complete bare and first-current values are
\begin{align}
 K_\Gamma&=-\frac12\langle\lambda_0,\lambda_1\rangle,\\
 K_\Gamma[01]_{+\dot a}
 &=\frac{\langle\lambda_0,\lambda_1\rangle}{6}
                       (-\lambda_{0,+\dot a}+\lambda_{1,+\dot a}),\\
 K_\Gamma[02]_{+\dot a}
 &=-\frac{\langle\lambda_0,\lambda_1\rangle}{6}
                       (2\lambda_{0,+\dot a}+\lambda_{1,+\dot a}),\\
 K_\Gamma[12]_{+\dot a}
 &=-\frac{\langle\lambda_0,\lambda_1\rangle}{6}
                       (\lambda_{0,+\dot a}+2\lambda_{1,+\dot a}).
 \label{networks:triangle-values}
\end{align}
For example, the entire order-$m$ tensor on edge $01$ is
\begin{equation}
 -\frac{\langle\lambda_0,\lambda_1\rangle}{(m+2)!}
 \sum_{S\subseteq\{1,\ldots,m\}}(-1)^{m-|S|}|S|!(m-|S|)!
 \prod_{j\in S}\lambda_{0,+\dot a_j}
 \prod_{j\notin S}\lambda_{1,+\dot a_j}.
 \label{networks:triangle-all-orders}
\end{equation}
This is the full tensor, with no free-index contraction.

For $\lambda_0=(1,0)$ and $\lambda_1=(0,1)$, the three complete cube
integrands combine to $-3/(1+x+y)^3$ on $[0,1]^2$.
Here $\tau_\Gamma=3$.  Since
\begin{equation}
 \int_0^1\!\int_0^1(1-x-y)^kdx\,dy
          =\frac{1+(-1)^k}{(k+1)(k+2)},
\end{equation}
the exact coefficients of \eqref{networks:binomial-series} are
\begin{equation}
 a_{2j}=-\frac3{2^{2j+3}},\qquad a_{2j+1}=0,\qquad
 \sum_{k=0}^{\infty}a_k=-\frac38\sum_{j=0}^{\infty}4^{-j}
                      =-\frac12.
 \label{networks:triangle-series}
\end{equation}
For requested width $1/1000$, the general bound first permits $K=18$:
\begin{equation}
 \sum_{k=0}^{18}a_k=-\frac{1048575}{2097152},\qquad
 B_{18}=\frac{87}{262144},\qquad
 \left|K_\Gamma-\sum_{k=0}^{18}a_k\right|=\frac1{2097152}.
 \label{networks:triangle-tail-control}
\end{equation}
The last equality is the sharper independently summed tail; the general
algorithm uses the middle bound.

\paragraph{A three-loop fork with zero bare value and all first currents.}
Order the edges as $(12,13,2A,2B,3A,3B,AB)$.  The sources are arbitrary
$\lambda_1,\lambda_A$, zero at $2,3$, and
$\lambda_B=-\lambda_1-\lambda_A$.  Two passive-star eliminations and the
remaining triangle give coordinates
\begin{equation}
 (r_1,r_2)\in(0,1)^2,\qquad
 (r_3,r_4)\in\Delta^2,\qquad(r_5,r_6)\in\Delta^2.
 \label{networks:fork-domain}
\end{equation}
To display every current without suppressing a source, write the three
retained triangle currents explicitly as
\begin{equation}
 y_{1A}^{\rm red}=r_5\lambda_1-r_6\lambda_A,\quad
 y_{1B}^{\rm red}=(1-r_5)\lambda_1+r_6\lambda_A,\quad
 y_{AB}^{\rm red}=r_5\lambda_1+(1-r_6)\lambda_A.
 \label{networks:fork-reduced}
\end{equation}
Each equation below is an equality of complete spinors:
\begin{align}
 y_{12}&=r_1y_{1A}^{\rm red}+r_2y_{1B}^{\rm red},&
 y_{13}&=(1-r_1)y_{1A}^{\rm red}+(1-r_2)y_{1B}^{\rm red},\\
 y_{2A}&=r_1y_{1A}^{\rm red}-r_3y_{AB}^{\rm red},&
 y_{2B}&=r_2y_{1B}^{\rm red}+r_3y_{AB}^{\rm red},\\
 y_{3A}&=(1-r_1)y_{1A}^{\rm red}-r_4y_{AB}^{\rm red},&
 y_{3B}&=(1-r_2)y_{1B}^{\rm red}+r_4y_{AB}^{\rm red},\\
 y_{AB}&=(1-r_3-r_4)y_{AB}^{\rm red}.&&
 \label{networks:fork-all-currents}
\end{align}
The full density, including orientation, is
\begin{equation}
 \Omega_\Gamma=\langle\lambda_1,\lambda_A\rangle^3
 (r_1-r_2)(1-r_3-r_4)r_5(1-r_5-r_6)
                       \,dr_1\cdots dr_6.
 \label{networks:fork-density}
\end{equation}
For the tree $(12,13,2A,2B)$, compute the full Jacobian of the chord
spinors $(y_{3A},y_{3B},y_{AB})$.  Its determinant is minus the scalar
polynomial in \eqref{networks:fork-density}.  The sign
$(-1)^{\binom32}=-1$ in \eqref{networks:current-area} cancels it.
The resistance chart preserves projective orientation: the two inverse
star changes reverse orientation twice, and their parallel fractions and
the final triangle have the displayed positive orientation.  The
independent replay also checks the full original resistance Jacobian and
original-network density at exact interior points.

The bare integral vanishes because $\int_0^1\int_0^1(r_1-r_2)dr_1dr_2=0$.
All seven complete first-current spinors are nevertheless determined.
Each row below gives $K_\Gamma[e]_{+\dot a}$ as
$\langle\lambda_1,\lambda_A\rangle^3/8640$ times the indicated spinor:
\begin{equation}
 \begin{array}{c|c}
 e & \text{complete spinor factor}\\
 \hline
 12&-\lambda_{1,+\dot a}-2\lambda_{A,+\dot a}\\
 13&\lambda_{1,+\dot a}+2\lambda_{A,+\dot a}\\
 2A&2\lambda_{1,+\dot a}-\lambda_{A,+\dot a}\\
 2B&-3\lambda_{1,+\dot a}-\lambda_{A,+\dot a}\\
 3A&-2\lambda_{1,+\dot a}+\lambda_{A,+\dot a}\\
 3B&3\lambda_{1,+\dot a}+\lambda_{A,+\dot a}\\
 AB&0
 \end{array}
 \label{networks:fork-values}
\end{equation}
The table denotes polynomial identities after multiplication by the
displayed bracket cubed; it remains valid when that bracket vanishes.
To calculate it directly, the two interval moments are $1/12$ and
$-1/12$ after multiplication by $r_1$ and $r_2$ respectively, while
\begin{align}
 \int_{\Delta^2}(1-r_3-r_4)dr_3dr_4&=\frac16,\\
 \int_{\Delta^2}r_5(1-r_5-r_6)y_{1A}^{\rm red}dr_5dr_6
                           &=\frac{2\lambda_1-\lambda_A}{120},\\
 \int_{\Delta^2}r_5(1-r_5-r_6)y_{1B}^{\rm red}dr_5dr_6
                           &=\frac{3\lambda_1+\lambda_A}{120}.
 \label{networks:fork-factorized-moments}
\end{align}
All terms independent of $r_1,r_2$ cancel before these remaining finite
integrals.  The full polynomial calculation independently gives the same
values and all integrated Kirchhoff identities.  Its executable and
certificate are listed in Appendix~\ref{app:certificates}.

\subsection{Relation to finite periods and physical assembly}
\label{networks:scope}

Theorem~\ref{networks:rational-series} evaluates every finite original
current word on every valid product forest by an exact infinite series
with finite specification and a finite order-to-error rule.  The chain
reduction \eqref{networks:chain-factor} can instead terminate in a finite
period expression when the required scalar values are known.  Theorem
\ref{periods:wheel-span} supplies such a uniform classification for wheel
cores.  The finite census and source-polynomial calculations in the
computational appendix have their separately stated topological, source,
and derivative ranges.

Finally, $K_\Gamma$ is an oriented geometric kernel.  It contains neither
the full Wick matching sum nor the color, representation, coupling,
Koszul, and propagator-normalization factors of a physical product.
Vanishing of a bare geometric kernel does not imply vanishing of its
current moments, as \eqref{networks:fork-values} shows.  Cancellation of
periods after complete physical assembly must be proved for the specified
fields and contractions; it is not a consequence of sourced circuit
elimination or of the convergent rational series.
